% Apartado 4 — Responsable: Apunte Pazmiño Washington German
\section{Relación con la Unidad 1}\label{sec:unidad1}

\textbf{Transparencia.} Se rompió la transparencia de fallos: según su definición, el usuario no debería notar el fallo y el sistema tendría que recuperarse solo~\cite{vansteen2016brief}. Pasó lo contrario: en algunas regiones no se podía entrar a sitios que usan Cloudflare, mientras que en otras todo funcionaba con normalidad. Tampoco hubo recuperación automática: los ingenieros tuvieron que revertir el cambio y todo volvió a la normalidad recién a las 07:42 UTC~\cite{cloudflare2022}.

\textbf{Teorema CAP.} Sí hubo una partición, pero provocada por una configuración y no por un cable cortado: al dejar de anunciarse los prefijos BGP \emph{site-local}, los servidores de los 19 centros ya no se alcanzaban entre sí ni llegaban a los servidores de origen de los clientes~\cite{cloudflare2022}. Ante una partición hay que elegir: responder con datos desactualizados (se pierde la consistencia) o no responder (se pierde la disponibilidad)~\cite{gilbert2002brewer}. Pasó lo segundo: se sacrificó la disponibilidad, aunque nadie lo eligió. El informe no menciona pérdida de datos~\cite{cloudflare2022}.

\textbf{Tolerancia a fallos.} Debía proteger el despliegue escalonado, junto con la prueba en seco y la revisión por pares, pero ninguno detectó el error: el escalonado no incluyó un centro MCP hasta el último paso, así que el cambio llegó de golpe a los \emph{spines} de los 19 centros~\cite{cloudflare2022}. La malla de \emph{spines} tampoco ayudó: la redundancia cubre fallos independientes, no una mala configuración copiada en todos. Esto coincide con Govindan \emph{et al.}: en la mayoría de los fallos graves que analizaron en la red de Google había una operación de gestión en curso~\cite{govindan2016evolve}.
